\numoptstrue
\yearpaper{2024}{5 May 2024 \textperiodcentered{} Test Booklet Code S6 \textperiodcentered{} Physics Q 1--50 (Section A: 1--35, Section B: 36--50, any 10 of 15) \textperiodcentered{} 180 Marks}{}
\begin{multicols}{2}

\QQ{Match List I with List II.\par\smallskip
\begin{tabular}{@{}lp{0.42\linewidth}lp{0.25\linewidth}@{}}
\multicolumn{2}{@{}p{0.5\linewidth}}{\textit{List I} (spectral lines of hydrogen for transitions from)} & \multicolumn{2}{p{0.34\linewidth}}{\textit{List II} (wavelength, nm)}\\
A. & $n_2 = 3$ to $n_1 = 2$ & I. & 410.2\\
B. & $n_2 = 4$ to $n_1 = 2$ & II. & 434.1\\
C. & $n_2 = 5$ to $n_1 = 2$ & III. & 656.3\\
D. & $n_2 = 6$ to $n_1 = 2$ & IV. & 486.1
\end{tabular}\par\smallskip
Choose the correct answer from the options given below:}
{\opt{A-I, B-II, C-III, D-IV}{A-II, B-I, C-IV, D-III}{A-III, B-IV, C-II, D-I}{A-IV, B-III, C-I, D-II}}

\QQ{Match List-I with List-II.\par\smallskip
\begin{tabular}{@{}ll@{\ }lp{0.34\linewidth}@{}}
\multicolumn{2}{@{}l}{\textit{List-I} (Material)} & \multicolumn{2}{l}{\textit{List-II} (Susceptibility $\chi$)}\\
A. & Diamagnetic & I. & $\chi = 0$\\
B. & Ferromagnetic & II. & $0>\chi\ge-1$\\
C. & Paramagnetic & III. & $\chi\gg1$\\
D. & Non-magnetic & IV. & $0<\chi<\varepsilon$ (a small positive number)
\end{tabular}\par\smallskip
Choose the correct answer from the options given below:}
{\opt{A-IV, B-III, C-II, D-I}{A-II, B-III, C-IV, D-I}{A-II, B-I, C-III, D-IV}{A-III, B-II, C-I, D-IV}}

\QQ{In a uniform magnetic field of $0.049\unit{T}$, a magnetic needle performs 20 complete oscillations in 5 seconds as shown. The moment of inertia of the needle is $9.8\times10^{-6}\unit{kg\,m^2}$. If the magnitude of magnetic moment of the needle is $x\times10^{-5}\unit{A\,m^2}$, then the value of $x$ is:
\fig{\begin{tikzpicture}[scale=0.8]
\fill[pattern=north east lines] (-0.6,1.2) rectangle (0.6,1.4); \draw (-0.6,1.2) -- (0.6,1.2); \draw (0,1.2) -- (0,0.15);
\draw[thick] (-0.7,-0.15) -- (0.7,0.15); \node[lbl] at (0.85,0.3) {N}; \node[lbl] at (-0.85,-0.3) {S};
\draw[densely dashed,->] (-1.2,0.4) -- (1.4,0.4); \draw[densely dashed,->] (-1.2,-0.4) -- (1.4,-0.4) node[right,lbl]{$B$};
\end{tikzpicture}}}
{\opt{$1280\pi^2$}{$5\pi^2$}{$128\pi^2$}{$50\pi^2$}}

\QQ{An unpolarised light beam strikes a glass surface at Brewster's angle. Then}
{\opt{the reflected light will be completely polarised but the refracted light will be partially polarised.}{the reflected light will be partially polarised.}{the refracted light will be completely polarised.}{both the reflected and refracted light will be completely polarised.}}

\QQ{Consider the following statements A and B and identify the correct answer:
\fig{\begin{tikzpicture}[scale=0.7]
\draw[->] (-1.2,0) -- (1.4,0) node[right,lbl]{$V$}; \draw[->] (0,-0.9) -- (0,1.2) node[above,lbl]{$I$};
\node[lbl] at (0.5,0.45) {(I)}; \node[lbl] at (-0.55,0.45) {(II)}; \node[lbl] at (-0.55,-0.45) {(III)}; \node[lbl] at (0.5,-0.45) {(IV)};
\end{tikzpicture}}
\begin{stmts}[\Alph*.]
\item For a solar-cell, the I-V characteristics lies in the IV quadrant of the given graph.
\item In a reverse biased pn junction diode, the current measured in ($\mu$A), is due to majority charge carriers.
\end{stmts}}
{\opt{Both A and B are incorrect.}{A is correct but B is incorrect.}{A is incorrect but B is correct.}{Both A and B are correct.}}

\QQ{A light ray enters through a right angled prism at point P with the angle of incidence $30^\circ$ as shown in figure. It travels through the prism parallel to its base BC and emerges along the face AC. The refractive index of the prism is:
\fig{\begin{tikzpicture}[scale=0.85]
\coordinate (A) at (0.6,2.0); \coordinate (B) at (0,0); \coordinate (C) at (3.0,0);
\draw (A) -- (B) -- (C) -- cycle; \node[lbl,above] at (A) {A}; \node[lbl,below left] at (B) {B}; \node[lbl,below right] at (C) {C};
\coordinate (P) at (0.3,1.0); \node[lbl,left] at (P) {P};
\draw[thick,->] (-0.7,1.3) -- (P); \draw[thick,->] (P) -- (1.5,1.0); \draw[thick] (1.5,1.0) -- (1.79,1.0);
\draw[thick,->] (1.79,1.0) -- (2.6,0.48);
\draw[densely dashed] (P) -- +(-0.8,0.24) (P) -- +(0.8,-0.24); \node[lbl] at (-0.35,0.9) {$30^\circ$};
\end{tikzpicture}}}
{\opt{$\dfrac{\sqrt3}2$}{$\dfrac{\sqrt5}4$}{$\dfrac{\sqrt5}2$}{$\dfrac{\sqrt3}4$}}

\QQ{A particle moving with uniform speed in a circular path maintains:}
{\opt{varying velocity and varying acceleration.}{constant velocity.}{constant acceleration.}{constant velocity but varying acceleration.}}

\QQ{The graph which shows the variation of $\left(\dfrac1{\lambda^2}\right)$ and its kinetic energy, $E$ is (where $\lambda$ is de Broglie wavelength of a free particle):}
{\optfig{\gaxes{$E$}{$\frac1{\lambda^2}$}{\draw (0,0) -- (1.6,1.5);}}
{\gaxes{$E$}{$\frac1{\lambda^2}$}{\draw (0,0) .. controls (1.0,0.1) and (1.4,0.6) .. (1.7,1.6);}}
{\gaxes{$E$}{$\frac1{\lambda^2}$}{\draw (0.1,1.6) .. controls (0.4,0.5) and (0.9,0.2) .. (1.8,0.1);}}
{\gaxes{$E$}{$\frac1{\lambda^2}$}{\draw (0,0) .. controls (0.2,1.0) and (0.7,1.5) .. (1.7,1.6);}}}

\QQ{A wire of length `$l$' and resistance $100\unit{\Omega}$ is divided into 10 equal parts. The first 5 parts are connected in series while the next 5 parts are connected in parallel. The two combinations are again connected in series. The resistance of this final combination is:}
{\opt{60\,$\Omega$}{26\,$\Omega$}{52\,$\Omega$}{55\,$\Omega$}}

\QQ{The moment of inertia of a thin rod about an axis passing through its mid point and perpendicular to the rod is $2400\unit{g\,cm^2}$. The length of the 400\,g rod is nearly:}
{\opt{72.0\,cm}{8.5\,cm}{17.5\,cm}{20.7\,cm}}

\QQ{The output ($Y$) of the given logic gate is similar to the output of an/a:
\fig{\begin{circuitikz}[scale=0.7,transform shape]
\draw (0,1.0) node[nand port](G1){}; \draw (0,-0.6) node[nor port](G2){}; \draw (2.5,0.2) node[nor port](G3){};
\draw (G1.in 1) -- ++(-0.3,0) coordinate(a1); \draw (G1.in 2) -- ++(-0.3,0) coordinate(a2); \draw (a1) -- (a2); \draw ($(a1)!0.5!(a2)$) -- ++(-0.4,0) node[left]{$A$};
\draw (G2.in 1) -- ++(-0.3,0) coordinate(b1); \draw (G2.in 2) -- ++(-0.3,0) coordinate(b2); \draw (b1) -- (b2); \draw ($(b1)!0.5!(b2)$) -- ++(-0.4,0) node[left]{$B$};
\draw (G1.out) -| (G3.in 1); \draw (G2.out) -| (G3.in 2); \draw (G3.out) -- ++(0.4,0) node[right]{$Y$};
\end{circuitikz}}}
{\opt{AND gate}{NAND gate}{NOR gate}{OR gate}}

\QQ{In a vernier calipers, $(N+1)$ divisions of vernier scale coincide with $N$ divisions of main scale. If 1 MSD represents $0.1\unit{mm}$, the vernier constant (in cm) is:}
{\opt{$10(N+1)$}{$\dfrac1{10N}$}{$\dfrac1{100(N+1)}$}{$100N$}}

\QQ{At any instant of time $t$, the displacement of any particle is given by $2t-1$ (SI unit) under the influence of force of $5\unit{N}$. The value of instantaneous power is (in SI unit):}
{\opt{6}{10}{5}{7}}

\QQ{A thin spherical shell is charged by some source. The potential difference between the two points C and P (in V) shown in the figure is: (Take $\dfrac1{4\pi\varepsilon_0} = 9\times10^9$ SI units)
\fig{\begin{tikzpicture}[scale=0.75]
\draw (0,0) circle (1); \foreach \a in {0,30,...,330} \node[font=\tiny] at (\a:1.12) {$+$};
\fill (0,0) circle (1.3pt) node[above left,lbl]{C}; \fill (25:1) circle (1.3pt) node[right=3pt,lbl]{P};
\draw[densely dotted] (0,0) -- (25:1); \draw (0,0) -- (-125:1) node[pos=0.55,right,lbl]{$R = 3\unit{cm}$};
\node[lbl] at (0,-1.5) {$q = 1\unit{\mu C}$};
\end{tikzpicture}}}
{\opt{zero}{$3\times10^5$}{$1\times10^5$}{$0.5\times10^5$}}

\QQ{A thin flat circular disc of radius $4.5\unit{cm}$ is placed gently over the surface of water. If surface tension of water is $0.07\unit{N\,m^{-1}}$, then the excess force required to take it away from the surface is:}
{\opt{99\,N}{19.8\,mN}{198\,N}{1.98\,mN}}

\QQ{A logic circuit provides the output $Y$ as per the following truth table:\par\smallskip
{\centering\begin{tabular}{|c|c|c|}\hline $A$ & $B$ & $Y$\\\hline 0 & 0 & 1\\ 0 & 1 & 0\\ 1 & 0 & 1\\ 1 & 1 & 0\\\hline\end{tabular}\par}\smallskip
The expression for the output $Y$ is:}
{\opt{$B$}{$A\cdot B+\bar A$}{$A\cdot\bar B+\bar A$}{$\bar B$}}

\QQ{Given below are two statements: one is labelled as \textbf{Assertion A} and the other is labelled as \textbf{Reason R}.\par
\textbf{Assertion A}: The potential ($V$) at any axial point, at $2\unit{m}$ distance ($r$) from the centre of the dipole of dipole moment vector $\vec P$ of magnitude $4\times10^{-6}\unit{C\,m}$, is $\pm9\times10^3\unit{V}$. (Take $\dfrac1{4\pi\varepsilon_0} = 9\times10^9$ SI units)\par
\textbf{Reason R}: $V = \pm\dfrac{2P}{4\pi\varepsilon_0r^2}$, where $r$ is the distance of any axial point, situated at $2\unit{m}$ from the centre of the dipole.\par
In the light of the above statements, choose the \textit{correct} answer from the options given below:}
{\opt{A is false but R is true.}{Both A and R are true and R is the correct explanation of A.}{Both A and R are true and R is NOT the correct explanation of A.}{A is true but R is false.}}

\QQ{The terminal voltage of the battery, whose emf is $10\unit{V}$ and internal resistance $1\unit{\Omega}$, when connected through an external resistance of $4\unit{\Omega}$ as shown in the figure is:
\fig{\begin{circuitikz}[scale=0.8,transform shape,font=\Large]
\draw (0,2.2) to[R,l=4\,$\Omega$] (4,2.2) -- (4,0.8);
\draw[->] (0,1.0) -- (0,1.7); \draw (0,0.8) -- (0,2.2); \draw[->] (4,1.7) -- (4,1.0);
\draw (0,0.8) -- (0.6,0.8) to[battery1,l_=10\,V] (1.8,0.8) to[R,l=1\,$\Omega$] (3.4,0.8) -- (4,0.8);
\draw[densely dashed] (0.4,0.2) rectangle (3.6,1.5);
\end{circuitikz}}}
{\opt{10\,V}{4\,V}{6\,V}{8\,V}}

\QQ{\fig{\begin{tikzpicture}[scale=1.0]
\foreach \x in {0,0.25,...,1.0} \draw (\x,0) .. controls (\x+0.1,0.6) and (\x+0.15,0.6) .. (\x+0.25,0);
\node[lbl,below] at (0,0) {A}; \node[lbl,below] at (1.25,0) {B}; \node[lbl] at (0.6,-0.65) {Solenoid-1};
\draw (1.9,0) rectangle (3.3,0.6); \node[lbl] at (2.2,0.3) {N}; \node[lbl] at (3.0,0.3) {S};
\draw[->] (3.4,0.3) -- (3.9,0.3) node[above,lbl]{$v$};
\begin{scope}[xshift=4.2cm]
\foreach \x in {0,0.25,...,1.0} \draw (\x,0) .. controls (\x+0.1,0.6) and (\x+0.15,0.6) .. (\x+0.25,0);
\node[lbl,below] at (0,0) {C}; \node[lbl,below] at (1.25,0) {D}; \node[lbl] at (0.6,-0.65) {Solenoid-2};
\end{scope}
\end{tikzpicture}}
In the above diagram, a strong bar magnet is moving towards solenoid-2 from solenoid-1. The direction of induced current in solenoid-1 and that in solenoid-2, respectively, are through the directions:}
{\opt{BA and DC}{AB and DC}{BA and CD}{AB and CD}}

\QQ{In an ideal transformer, the turns ratio is $\dfrac{N_p}{N_s} = \dfrac12$. The ratio $V_s:V_p$ is equal to (the symbols carry their usual meaning):}
{\opt{1:4}{1:2}{2:1}{1:1}}

\QQ{A tightly wound 100 turns coil of radius $10\unit{cm}$ carries a current of $7\unit{A}$. The magnitude of the magnetic field at the centre of the coil is (Take permeability of free space as $4\pi\times10^{-7}$ SI units):}
{\opt{44\,T}{44\,mT}{4.4\,T}{4.4\,mT}}

\QQ{Given below are two statements:\par
\textbf{Statement I}: Atoms are electrically neutral as they contain equal number of positive and negative charges.\par
\textbf{Statement II}: Atoms of each element are stable and emit their characteristic spectrum.\par
In the light of the above statements, choose the \textit{most appropriate} answer from the options given below:}
{\opt{Statement I is incorrect but Statement II is correct.}{Both Statement I and Statement II are correct.}{Both Statement I and Statement II are incorrect.}{Statement I is correct but Statement II is incorrect.}}

\QQ{The quantities which have the same dimensions as those of solid angle are:}
{\opt{angular speed and stress}{strain and angle}{stress and angle}{strain and arc}}

\QQ{If $c$ is the velocity of light in free space, the correct statements about photon among the following are:
\begin{stmts}[\Alph*.]
\item The energy of a photon is $E = h\nu$.
\item The velocity of a photon is $c$.
\item The momentum of a photon, $p = \dfrac{h\nu}c$.
\item In a photon-electron collision, both total energy and total momentum are conserved.
\item Photon possesses positive charge.
\end{stmts}
Choose the correct answer from the options given below:}
{\opt{A, B, D and E only}{A and B only}{A, B, C and D only}{A, C and D only}}

\QQ{A bob is whirled in a horizontal plane by means of a string with an initial speed of $\omega$ rpm. The tension in the string is $T$. If speed becomes $2\omega$ while keeping the same radius, the tension in the string becomes:}
{\opt{$\sqrt2T$}{$T$}{$4T$}{$\dfrac T4$}}

\QQ{If $x = 5\sin\left(\pi t+\dfrac\pi3\right)\unit{m}$ represents the motion of a particle executing simple harmonic motion, the amplitude and time period of motion, respectively, are:}
{\opt{5\,m, 1\,s}{5\,cm, 2\,s}{5\,m, 2\,s}{5\,cm, 1\,s}}

\QQ{A thermodynamic system is taken through the cycle $abcda$. The work done by the gas along the path $bc$ is:
\fig{\begin{tikzpicture}[xscale=0.006,yscale=0.006]
\draw[->] (0,0) -- (520,0) node[right,lbl]{$V$}; \draw[->] (0,0) -- (0,370) node[above,lbl]{$P$};
\draw[densely dashed] (0,300) node[left,lbl]{300\,kPa} -- (100,300) (0,100) node[left,lbl]{100\,kPa} -- (100,100) (100,100) -- (100,0) node[below,lbl]{100\,cm$^3$} (400,100) -- (400,0) node[below,lbl]{400\,cm$^3$};
\draw[thick,postaction={decorate},decoration={markings,mark=at position 0.15 with {\arrow{Stealth}},mark=at position 0.42 with {\arrow{Stealth}},mark=at position 0.65 with {\arrow{Stealth}},mark=at position 0.92 with {\arrow{Stealth}}}] (100,100) -- (400,100) -- (400,300) -- (100,300) -- cycle;
\node[lbl,below right] at (100,100) {$a$}; \node[lbl,right] at (400,100) {$b$}; \node[lbl,above right] at (400,300) {$c$}; \node[lbl,above right] at (100,300) {$d$};
\end{tikzpicture}}}
{\opt{$-60\unit{J}$}{zero}{30\,J}{$-90\unit{J}$}}

\QQ{${}^{290}_{\ 82}X\xrightarrow{\ \alpha\ }Y\xrightarrow{\ e^+\ }Z\xrightarrow{\ \beta^-\ }P\xrightarrow{\ e^-\ }Q$. In the nuclear emission stated above, the mass number and atomic number of the product Q respectively, are:}
{\opt{286, 81}{280, 81}{286, 80}{288, 82}}

\QQ{Two bodies A and B of same mass undergo completely inelastic one dimensional collision. The body A moves with velocity $v_1$ while body B is at rest before collision. The velocity of the system after collision is $v_2$. The ratio $v_1:v_2$ is:}
{\opt{1:4}{1:2}{2:1}{4:1}}

\QQ{In the following circuit, the equivalent capacitance between terminal A and terminal B is:
\fig{\begin{circuitikz}[scale=0.6,transform shape,font=\Large]
\draw (0,1.2) node[ocirc]{} node[left]{A} -- (1,1.2) -- (1,2.4) to[C,l=2\,$\mu$F] (3,2.4) to[C,l=2\,$\mu$F] (5,2.4) -- (5,1.2) -- (6,1.2) node[ocirc]{} node[right]{B};
\draw (1,1.2) -- (1,0) to[C,l_=2\,$\mu$F] (3,0) to[C,l_=2\,$\mu$F] (5,0) -- (5,1.2);
\draw (3,2.4) to[C,l=2\,$\mu$F] (3,0);
\end{circuitikz}}}
{\opt{4\,$\mu$F}{2\,$\mu$F}{1\,$\mu$F}{0.5\,$\mu$F}}

\QQ{A horizontal force $10\unit{N}$ is applied to a block A as shown in figure. The mass of blocks A and B are $2\unit{kg}$ and $3\unit{kg}$, respectively. The blocks slide over a frictionless surface. The force exerted by block A on block B is:
\fig{\begin{tikzpicture}[scale=0.7]
\fill[pattern=north east lines] (-0.4,-0.2) rectangle (3.6,0); \draw (-0.4,0) -- (3.6,0);
\draw (1,0) rectangle (2,0.8) node[midway,lbl,align=center]{A\\2\,kg}; \draw (2,0) rectangle (3.1,1.3) node[midway,lbl,align=center]{B\\3\,kg};
\draw[->] (-0.3,0.4) -- (1,0.4) node[pos=0.4,above,lbl]{$F = 10\unit{N}$};
\end{tikzpicture}}}
{\opt{10\,N}{zero}{4\,N}{6\,N}}

\QQ{A wheel of a bullock cart is rolling on a level road as shown in the figure below. If its linear speed is $v$ in the direction shown, which one of the following options is correct (P and Q are any highest and lowest points on the wheel, respectively)?
\fig{\begin{tikzpicture}[scale=0.6]
\fill[pattern=north east lines] (-1.6,-1.2) rectangle (1.8,-1.0); \draw (-1.6,-1) -- (1.8,-1);
\draw (0,0) circle (1); \foreach \a in {0,45,...,315} \draw (0,0) -- (\a:1);
\node[lbl,above] at (0,1) {P}; \node[lbl,below] at (0,-1.2) {Q}; \draw[->] (1.3,0) -- (2.2,0) node[above,lbl]{$v$};
\end{tikzpicture}}}
{\opt{Point P has zero speed.}{Point P moves slower than point Q.}{Point P moves faster than point Q.}{Both the points P and Q move with equal speed.}}

\QQ{The mass of a planet is $\dfrac1{10}$th that of the earth and its diameter is half that of the earth. The acceleration due to gravity on that planet is:}
{\opt{$3.92\unit{m\,s^{-2}}$}{$19.6\unit{m\,s^{-2}}$}{$9.8\unit{m\,s^{-2}}$}{$4.9\unit{m\,s^{-2}}$}}

\QQ{The maximum elongation of a steel wire of $1\unit{m}$ length if the elastic limit of steel and its Young's modulus, respectively, are $8\times10^8\unit{N\,m^{-2}}$ and $2\times10^{11}\unit{N\,m^{-2}}$, is:}
{\opt{8\,mm}{4\,mm}{0.4\,mm}{40\,mm}}

\QQ{If the monochromatic source in Young's double slit experiment is replaced by white light, then}
{\opt{all bright fringes will be of equal width.}{interference pattern will disappear.}{there will be a central dark fringe surrounded by a few coloured fringes.}{there will be a central bright white fringe surrounded by a few coloured fringes.}}

\QQ{A small telescope has an objective of focal length $140\unit{cm}$ and an eye piece of focal length $5.0\unit{cm}$. The magnifying power of telescope for viewing a distant object is:}
{\opt{32}{34}{28}{17}}

\QQ{If the plates of a parallel plate capacitor connected to a battery are moved close to each other, then
\begin{stmts}[\Alph*.]
\item the charge stored in it, increases.
\item the energy stored in it, decreases.
\item its capacitance increases.
\item the ratio of charge to its potential remains the same.
\item the product of charge and voltage increases.
\end{stmts}
Choose the most appropriate answer from the options given below:}
{\opt{A, B and C only}{A, B and E only}{A, C and E only}{B, D and E only}}

\QQ{A $10\unit{\mu F}$ capacitor is connected to a $210\unit{V}$, $50\unit{Hz}$ source as shown in figure. The peak current in the circuit is nearly ($\pi = 3.14$):
\fig{\begin{circuitikz}[scale=0.6,transform shape,font=\Large]
\draw (0,0) to[sV] (0,1.8) -- (0,2.4) to[C,l={$C = 10\unit{\mu F}$}] (3,2.4) -- (3,0) -- (0,0); \node at (1.5,-0.5) {210\,V, 50\,Hz};
\end{circuitikz}}}
{\opt{0.35\,A}{0.58\,A}{0.93\,A}{1.20\,A}}

\QQ{Choose the correct circuit which can achieve the bridge balance. (Each figure: arms $10\,\Omega$ and $10\,\Omega$ above, $15\,\Omega$ and $5\,\Omega$ below, galvanometer G across the bridge, diode D in the $5\,\Omega$ arm, cell E and key K.)}
{\optfig{\begin{circuitikz}[scale=0.42,transform shape]
\draw (0,0) to[R,l=10\,$\Omega$] (1.5,1.5) to[R,l=10\,$\Omega$] (3,0); \draw (0,0) to[R,l_=15\,$\Omega$] (1.5,-1.5) to[empty diode,l_=D] (2.25,-0.75) to[R,l_=5\,$\Omega$] (3,0);
\draw (1.5,1.5) to[rmeter,t=G] (1.5,-1.5); \draw (0,0) -- (-0.8,0) -- (-0.8,-2.6) to[battery1,l_=E] (1.5,-2.6) -- (3.8,-2.6) -- (3.8,0) -- (3,0);
\end{circuitikz}}
{\begin{circuitikz}[scale=0.42,transform shape]
\draw (0,0) to[R,l=10\,$\Omega$] (1.5,1.5) to[R,l=10\,$\Omega$] (3,0); \draw (0,0) to[R,l_=15\,$\Omega$] (1.5,-1.5) to[R,l_=5\,$\Omega$] (2.25,-0.75) to[empty diode,l_=D,invert] (3,0);
\draw (1.5,1.5) to[rmeter,t=G] (1.5,-1.5); \draw (0,0) -- (-0.8,0) -- (-0.8,-2.6) to[battery1,l_=E] (1.5,-2.6) -- (3.8,-2.6) -- (3.8,0) -- (3,0);
\end{circuitikz}}
{\begin{circuitikz}[scale=0.42,transform shape]
\draw (0,0) to[R,l=10\,$\Omega$] (1.5,1.5) to[R,l=10\,$\Omega$] (3,0); \draw (0,0) to[R,l_=15\,$\Omega$] (1.5,-1.5) to[R,l_=5\,$\Omega$] (2.25,-0.75) to[empty diode,l_=D] (3,0);
\draw (1.5,1.5) to[rmeter,t=G] (1.5,-1.5); \draw (0,0) -- (-0.8,0) -- (-0.8,-2.6) to[battery1,l_=E] (1.5,-2.6) -- (3.8,-2.6) -- (3.8,0) -- (3,0);
\end{circuitikz}}
{\begin{circuitikz}[scale=0.42,transform shape]
\draw (0,0) to[R,l=10\,$\Omega$] (1.5,1.5) to[R,l=10\,$\Omega$] (3,0); \draw (0,0) to[R,l_=15\,$\Omega$] (1.5,-1.5) to[empty diode,l=D] (3,0);
\draw (1.5,-1.5) -- (2.3,-1.9) to[R,l_=5\,$\Omega$] (3.4,-0.4) -- (3,0);
\draw (1.5,1.5) to[rmeter,t=G] (1.5,-1.5); \draw (0,0) -- (-0.8,0) -- (-0.8,-2.6) to[battery1,l_=E] (1.5,-2.6) -- (3.8,-2.6) -- (3.8,0) -- (3,0);
\end{circuitikz}}}

\QQ{Two heaters A and B have power rating of $1\unit{kW}$ and $2\unit{kW}$, respectively. Those two are first connected in series and then in parallel to a fixed power source. The ratio of power outputs for these two cases is:}
{\opt{2:3}{1:1}{2:9}{1:2}}

\QQ{The velocity ($v$)--time ($t$) plot of the motion of a body is shown below:
\fig{\gaxes{$t$ (s)}{$v$ (m\,s$^{-1}$)}{\draw (0.2,0) -- (0.6,1.1) -- (1.2,1.1) -- (1.6,0);}}
The acceleration ($a$)--time ($t$) graph that best suits this motion is:}
{\optfig{\gaxesm{$t$}{$a$}{\draw (0,0) -- (0.4,0) -- (0.4,-0.5) -- (1.0,-0.5) -- (1.0,0) -- (1.3,0) -- (1.9,0.6);}}
{\gaxesm{$t$}{$a$}{\draw (0,0) -- (0.5,0) -- (0.8,0.6) -- (1.3,0.6) -- (1.6,0) -- (2.0,0);}}
{\gaxesm{$t$}{$a$}{\draw (0,0) -- (0.3,0) -- (1.0,0.9) -- (1.9,0.9);}}
{\gaxesm{$t$}{$a$}{\draw (0,0) -- (0.4,0) -- (0.4,0.6) -- (0.8,0.6) -- (0.8,0) -- (1.4,0) -- (1.4,-0.6) -- (1.8,-0.6) -- (1.8,0);}}}

\QQ{If the mass of the bob in a simple pendulum is increased to thrice its original mass and its length is made half its original length, then the new time period of oscillation is $\dfrac x2$ times its original time period. Then the value of $x$ is:}
{\opt{4}{$\sqrt3$}{$\sqrt2$}{$2\sqrt3$}}

\QQ{A metallic bar of Young's modulus, $0.5\times10^{11}\unit{N\,m^{-2}}$ and coefficient of linear thermal expansion $10^{-5}\,{}^\circ\mathrm C^{-1}$, length $1\unit{m}$ and area of cross-section $10^{-3}\unit{m^2}$ is heated from $0^\circ$C to $100^\circ$C without expansion or bending. The compressive force developed in it is:}
{\opt{$2\times10^3\unit{N}$}{$5\times10^3\unit{N}$}{$50\times10^3\unit{N}$}{$100\times10^3\unit{N}$}}

\QQ{A sheet is placed on a horizontal surface in front of a strong magnetic pole. A force is needed to:
\begin{stmts}[\Alph*.]
\item hold the sheet there if it is magnetic.
\item hold the sheet there if it is non-magnetic.
\item move the sheet away from the pole with uniform velocity if it is conducting.
\item move the sheet away from the pole with uniform velocity if it is both, non-conducting and non-polar.
\end{stmts}
Choose the correct statement(s) from the options given below:}
{\opt{C only}{B and D only}{A and C only}{A, C and D only}}

\QQ{The property which is not of an electromagnetic wave travelling in free space is that:}
{\opt{they originate from charges moving with uniform speed.}{they are transverse in nature.}{the energy density in electric field is equal to energy density in magnetic field.}{they travel with a speed equal to $\dfrac1{\sqrt{\mu_0\varepsilon_0}}$.}}

\QQ{A parallel plate capacitor is charged by connecting it to a battery through a resistor. If $I$ is the current in the circuit, then in the gap between the plates:}
{\opt{displacement current of magnitude greater than $I$ flows but can be in any direction.}{there is no current.}{displacement current of magnitude equal to $I$ flows in the same direction as $I$.}{displacement current of magnitude equal to $I$ flows in a direction opposite to that of $I$.}}

\QQ{An iron bar of length L has magnetic moment M. It is bent at the middle of its length such that the two arms make an angle $60^\circ$ with each other. The magnetic moment of this new magnet is:}
{\opt{$\dfrac M{\sqrt3}$}{M}{$\dfrac M2$}{2M}}

\QQ{The minimum energy required to launch a satellite of mass $m$ from the surface of earth of mass $M$ and radius $R$ in a circular orbit at an altitude of $2R$ from the surface of the earth is:}
{\opt{$\dfrac{GmM}{3R}$}{$\dfrac{5GmM}{6R}$}{$\dfrac{2GmM}{3R}$}{$\dfrac{GmM}{2R}$}}

\QQ{The following graph represents the T-V curves of an ideal gas (where T is the temperature and V the volume) at three pressures $P_1$, $P_2$ and $P_3$ compared with those of Charles's law represented as dotted lines.
\fig{\gaxes{$V$}{$T$}{\draw (0.15,0.1) .. controls (0.4,0.7) and (0.6,1.2) .. (0.7,1.7) node[above,font=\tiny]{$P_1$}; \draw (0.25,0.1) .. controls (0.6,0.6) and (0.8,1.0) .. (1.0,1.5) node[above,font=\tiny]{$P_2$}; \draw (0.35,0.1) .. controls (0.8,0.45) and (1.1,0.8) .. (1.35,1.2) node[right,font=\tiny]{$P_3$}; \draw[densely dotted] (0,0) -- (0.65,1.7) (0,0) -- (0.95,1.5) (0,0) -- (1.35,1.2);}}
Then the correct relation is:}
{\opt{$P_1>P_2>P_3$}{$P_3>P_2>P_1$}{$P_1>P_3>P_2$}{$P_2>P_1>P_3$}}

\QQ{A force defined by $F = \alpha t^2+\beta t$ acts on a particle at a given time $t$. The factor which is dimensionless, if $\alpha$ and $\beta$ are constants, is:}
{\opt{$\alpha\beta/t$}{$\beta t/\alpha$}{$\alpha t/\beta$}{$\alpha\beta t$}}

\end{multicols}
